\section{Week 4: Qubits, Unitaries, Entanglement, Rank}\label{sec:week4}

Lectures of September 15 and 17.

\begin{opentable}
    \openrow{ex:basis-0915}{Exercise: entries of the column vector $\ket{00}$}
    \openrow{thm:unitary-0915}{Exercise: norm preserving iff unitary}
    \openrow{lem:hh-0915}{$HH = I$}
    \openrow{lem:cnot-0915}{Matrix of CNOT}
    \openrow{lem:xor-0915}{XOR is not reversible}
    \openrow{thm:bell-0915}{Exercise: the Bell state is entangled}
    \openrow{ex:example-rank-0917}{Exercise: rank of the example matrix}
    \openrow{thm:rank-equal-0917}{Exercise: row rank equals column rank}
    \openrow{thm:orth-ind-0917}{Exercise: orthogonal implies independent}
\end{opentable}

\subsection{Basis Vectors, Unitaries, Gates}

\begin{exercisebox}
\begin{pproblem}[Which entry is $1$?]\label{wp:basis}
    Write $\ket{10}$ as a column vector in $\mathbb{C}^4$, and say which entry of $\ket{101}$ in
    $\mathbb{C}^8$ is $1$.
\end{pproblem}
\tcblower
\begin{wsolution}
    \uses{\scr{ex:basis-0915}}
    The basis states are ordered by the binary number they spell: $\ket{00}, \ket{01}, \ket{10},
    \ket{11}$ are the entries $1, 2, 3, 4$. So
    $\ket{10} = \ket{1} \otimes \ket{0} = (0, 1)^T \otimes (1, 0)^T = (0, 0, 1, 0)^T$. The string
    $101$ is the number $5$, so $\ket{101}$ has a $1$ in entry $5 + 1 = 6$ and zeros elsewhere.
\end{wsolution}
\end{exercisebox}

\begin{exercisebox}
\begin{pproblem}[Is the matrix unitary?]\label{pp:unitary}
    Decide for each matrix whether it is unitary.
    \[U = \frac{1}{\sqrt{2}} \begin{pmatrix} 1 & i \\ i & 1 \end{pmatrix},
        \qquad
        V = \begin{pmatrix} 1 & 1 \\ 0 & 1 \end{pmatrix},
        \qquad
        W = \frac{1}{\sqrt{2}} \begin{pmatrix} 1 & 1 \\ 1 & 1 \end{pmatrix}.\]
\end{pproblem}
\tcblower
\begin{wsolution}
    \uses{\scr{thm:unitary-0915}}
    Compute $U^\dagger U$, where $\dagger$ means: transpose, and conjugate every entry.
    \[U^\dagger U = \frac{1}{2} \begin{pmatrix} 1 & -i \\ -i & 1 \end{pmatrix} \begin{pmatrix} 1 & i \\ i & 1 \end{pmatrix}
        = \frac{1}{2} \begin{pmatrix} 1 - i^2 & i - i \\ -i + i & -i^2 + 1 \end{pmatrix}
        = \begin{pmatrix} 1 & 0 \\ 0 & 1 \end{pmatrix}.\]
    So $U$ is unitary. Do not forget the conjugate: without it the diagonal would be $1 + i^2 = 0$.

    $V$ is not unitary: $V\ket{1} = \ket{0} + \ket{1}$ has norm $\sqrt{2}$, and a unitary matrix
    preserves norms (Sep 15). $W$ is not unitary: $W\ket{0} = W\ket{1} = \ket{+}$, so two orthogonal
    states are sent to the same state; also $W\left(\ket{0} - \ket{1}\right) = 0$.

    \emph{Quick test.} A matrix is unitary exactly when its columns are orthonormal, with the
    complex inner product.
\end{wsolution}
\end{exercisebox}


\begin{exercisebox}
\begin{pproblem}[$HXH$]\label{wp:hxh}
    Compute $HXH$ and $HZH$, where $X = \begin{pmatrix} 0 & 1 \\ 1 & 0 \end{pmatrix}$ and
    $Z = \begin{pmatrix} 1 & 0 \\ 0 & -1 \end{pmatrix}$.
\end{pproblem}
\tcblower
\begin{wsolution}
    \uses{\scr{lem:hh-0915}}
    $XH = \frac{1}{\sqrt{2}} \begin{pmatrix} 1 & -1 \\ 1 & 1 \end{pmatrix}$, so
    \[HXH = \frac{1}{2} \begin{pmatrix} 1 & 1 \\ 1 & -1 \end{pmatrix} \begin{pmatrix} 1 & -1 \\ 1 & 1 \end{pmatrix}
        = \frac{1}{2} \begin{pmatrix} 2 & 0 \\ 0 & -2 \end{pmatrix} = Z.\]
    Multiply $HXH = Z$ by $H$ on both sides and use $HH = I$: $X = HZH$. So $H$ swaps the bit flip
    and the phase flip. On states: $X$ fixes $\ket{+}$ and negates $\ket{-}$, just as $Z$ fixes
    $\ket{0}$ and negates $\ket{1}$.
\end{wsolution}
\end{exercisebox}

\subsection{Entangled or Product}

Write a two-qubit state as $a\ket{00} + b\ket{01} + c\ket{10} + d\ket{11}$. A product
$\left(\alpha_0\ket{0} + \alpha_1\ket{1}\right)\left(\beta_0\ket{0} + \beta_1\ket{1}\right)$ has
$a = \alpha_0\beta_0$, $b = \alpha_0\beta_1$, $c = \alpha_1\beta_0$, $d = \alpha_1\beta_1$, so
$ad = \alpha_0\alpha_1\beta_0\beta_1 = bc$. If $ad \neq bc$ the state is entangled; this is the Bell
state argument written once for all states. To show that a state is a product, write down the two
factors.

\begin{exercisebox}
\begin{pproblem}[Five states]\label{pp:five}
    Which of these states are entangled? Prove each answer.
    \begin{enumerate}[label=(\alph*)]
        \item $\frac{1}{2}\left(\ket{00} + \ket{01} + \ket{10} + \ket{11}\right)$
        \item $\frac{1}{2}\left(\ket{00} - \ket{01} + \ket{10} - \ket{11}\right)$
        \item $\frac{1}{2}\left(\ket{00} + \ket{01} + \ket{10} - \ket{11}\right)$
        \item $\frac{1}{\sqrt{2}}\left(\ket{01} - \ket{10}\right)$
        \item $\sqrt{\sfrac{1}{3}}\, \ket{00} + \sqrt{\sfrac{2}{3}}\, \ket{01}$
    \end{enumerate}
\end{pproblem}
\tcblower
\begin{wsolution}
    \uses{\scr{thm:bell-0915}}
    Compare $ad$ with $bc$:
    \begin{center}
        \begin{tabular}{cccl}
            \toprule
            State & $ad$              & $bc$              & Answer                                                                      \\ \midrule
            (a)   & $\sfrac{1}{4}$    & $\sfrac{1}{4}$    & Product: $\ket{+}\ket{+}$                                                   \\
            (b)   & $-\sfrac{1}{4}$   & $-\sfrac{1}{4}$   & Product: $\ket{+}\ket{-}$                                                   \\
            (c)   & $-\sfrac{1}{4}$   & $\sfrac{1}{4}$    & Entangled                                                                   \\
            (d)   & $0$               & $-\sfrac{1}{2}$   & Entangled                                                                   \\
            (e)   & $0$               & $0$               & Product: $\ket{0}\left(\sqrt{\sfrac{1}{3}}\,\ket{0} + \sqrt{\sfrac{2}{3}}\,\ket{1}\right)$ \\ \bottomrule
        \end{tabular}
    \end{center}

    \emph{Full write-up for (c).} Suppose the state equals
    $\left(\alpha_0\ket{0} + \alpha_1\ket{1}\right)\left(\beta_0\ket{0} + \beta_1\ket{1}\right)$.
    Then $\alpha_0\beta_0 = \frac{1}{2}$, $\alpha_0\beta_1 = \frac{1}{2}$,
    $\alpha_1\beta_0 = \frac{1}{2}$ and $\alpha_1\beta_1 = -\frac{1}{2}$. The product of the first
    and the last is $\alpha_0\alpha_1\beta_0\beta_1 = -\frac{1}{4}$. The product of the middle two
    is $\alpha_0\alpha_1\beta_0\beta_1 = +\frac{1}{4}$. This is a contradiction, so the state is
    entangled.
\end{wsolution}
\end{exercisebox}

\begin{exercisebox}
\begin{pproblem}[A small circuit]\label{pp:circuit}
    Apply $H$ to the first qubit and then CNOT (the first qubit is the control) to $\ket{10}$. Is
    the result entangled?
\end{pproblem}
\tcblower
\begin{wsolution}
    \uses{\scr{lem:cnot-0915}; \scr{thm:bell-0915}}
    \leavevmode
    \begin{enumerate}
        \item $H$ on the first qubit: $\ket{10}$ becomes
              $\frac{1}{\sqrt{2}}\left(\ket{0} - \ket{1}\right) \otimes \ket{0}
                  = \frac{1}{\sqrt{2}}\left(\ket{00} - \ket{10}\right)$.
        \item CNOT flips the second qubit when the first is $1$:
              $\frac{1}{\sqrt{2}}\left(\ket{00} - \ket{11}\right)$.
        \item Test: $ad = -\frac{1}{2}$ and $bc = 0$, so the state is entangled.
    \end{enumerate}
    The same circuit sends $\ket{00}$ to $\frac{1}{\sqrt{2}}\left(\ket{00} + \ket{11}\right)$,
    $\ket{01}$ to $\frac{1}{\sqrt{2}}\left(\ket{01} + \ket{10}\right)$ and $\ket{11}$ to
    $\frac{1}{\sqrt{2}}\left(\ket{01} - \ket{10}\right)$.
\end{wsolution}
\end{exercisebox}


\subsection{Rank and Independence}

\begin{exercisebox}
\begin{pproblem}[Rank of a $3 \times 3$ matrix]\label{wp:rank}
    Let
    \[A = \begin{pmatrix} 1 & 2 & 3 \\ 2 & 4 & 6 \\ 1 & 0 & 1 \end{pmatrix}.\]
    Find the column rank, a basis of $\operatorname{Im}(A)$, and the row rank, and write every row
    as a combination of two rows.
\end{pproblem}
\tcblower
\begin{wsolution}
    \uses{\scr{ex:example-rank-0917}; \scr{thm:rank-equal-0917}}
    \emph{Columns.} $c_1 = (1, 2, 1)^T$, $c_2 = (2, 4, 0)^T$, $c_3 = (3, 6, 1)^T = c_1 + c_2$. The
    columns $c_1$ and $c_2$ are independent: if $c_2 = \lambda c_1$, the third entries give
    $\lambda = 0$, but $c_2 \neq 0$. So $\operatorname{rank}_c(A) = 2$ and $\operatorname{Im}(A)
    = \operatorname{span}(c_1, c_2)$.

    \emph{Rows.} $r_1 = [1\ 2\ 3]$, $r_2 = [2\ 4\ 6]$, $r_3 = [1\ 0\ 1]$. Here $r_2 = 2r_1$, and
    $r_1$, $r_3$ are independent. So $\operatorname{rank}_r(A) = 2$, equal to the column rank as the
    exercise says. Every row in terms of $r_1, r_3$: $r_1 = 1 \cdot r_1 + 0 \cdot r_3$,
    $r_2 = 2 r_1 + 0 \cdot r_3$, $r_3 = 0 \cdot r_1 + 1 \cdot r_3$.
\end{wsolution}
\end{exercisebox}

\begin{exercisebox}
\begin{pproblem}[Orthogonal does not imply independent over $\set{0,1}$]\label{wp:gf2}
    Over $\set{0,1}^n$ with $x \cdot y = \sum_i x_i y_i \bmod 2$, is it still true that a nonzero
    $v$ orthogonal to independent $v_1, \dots, v_k$ is independent of them? Which step of the
    proof breaks?
\end{pproblem}
\tcblower
\begin{wsolution}
    \uses{\scr{thm:orth-ind-0917}}
    No. In $\set{0,1}^2$ take $v_1 = 11$ and $v = 11$. Then $v \cdot v_1 = 1 + 1 = 0 \bmod 2$, so $v$
    is nonzero and orthogonal to $v_1$, but $v = v_1$ is not independent of it.

    The proof takes $c\,v + \sum_i c_i v_i = 0$, takes the inner product with $v$ to get
    $c \norm{v}^2 = 0$, and concludes $c = 0$ from $\norm{v}^2 > 0$. Over $\set{0,1}$ the analogue
    of $\norm{v}^2$ is $v \cdot v$, the number of $1$s of $v$ mod $2$. It is $0$ for every $v$ with
    an even number of $1$s, so the step $c = 0$ fails.

    \emph{Looking ahead to week 6.} In Simon's algorithm every measured $y$ satisfies $s \cdot y = 0$, and $s$ itself may be one of
    them (when $s$ has an even number of $1$s, $s \cdot s = 0$). So independence of the $y$'s cannot
    come from orthogonality; it is a probability statement, proved separately in week~6
    (\scr{thm:indep-1001}).
\end{wsolution}
\end{exercisebox}
